\begin{frame}[allowframebreaks]{Pitfalls You Must Know Before Using Either}

    \begin{itemize}
        \setlength\itemsep{0.5em}
        \item \textbf{Correlated features.} Both methods perturb features independently, producing
              rows that could not exist. Attribution then gets spread arbitrarily among correlated
              features, and can be assigned to a feature the model barely uses. Group correlated
              features and attribute to the group.

        \item \textbf{The background / neighbourhood is a choice.} SHAP values are relative to the
              background dataset; LIME weights depend on the kernel width and the perturbation
              scheme. Report what you used --- otherwise the numbers are not reproducible.

        \item \textbf{Instability.} Re-running LIME with a different seed can reorder the top
              features. Always run it several times and check the explanation is stable before
              you show it to anyone.

        \item \textbf{Explanations can be attacked.} Slack et al.\ (2020) built classifiers that
              behave in a biased way on real data but detect the synthetic perturbation
              distributions used by LIME and SHAP, and answer innocuously there --- so the
              explanation hides the bias.
    \end{itemize}

    \vspace{0.5em}
    \begin{block}{}
        \centering
        \small \textbf{Optional lab:} run SHAP and LIME on the same tabular model, and watch LIME's top features
        change from one run to the next.
    \end{block}

    \vspace{0.3em}
    {\scriptsize D.~Slack, S.~Hilgard, E.~Jia, S.~Singh and H.~Lakkaraju, ``Fooling LIME and SHAP:
    Adversarial Attacks on Post hoc Explanation Methods,'' AIES 2020,
    \href{https://arxiv.org/abs/1911.02508v2}{arXiv:1911.02508v2}.}
\end{frame}

\begin{frame}{Explaining the Model vs Explaining the World}

    \begin{columns}[T]
        \begin{column}{0.48\textwidth}
            \begin{tcolorbox}[colback=raiBlue!5, colframe=raiBlue!60, title=What SHAP answers]
                \scriptsize ``How does \emph{this fitted function} distribute its output across the
                inputs I gave it?''

                \vspace{0.12em}
                A statement about the model. Correct as far as it goes.
            \end{tcolorbox}
        \end{column}
        \begin{column}{0.48\textwidth}
            \begin{tcolorbox}[colback=raiRed!5, colframe=raiRed!70, title=What people hear]
                \scriptsize ``Which factors \emph{cause} this outcome, and what should this person
                change?''

                \vspace{0.12em}
                A causal statement. Not supported by the attribution.
            \end{tcolorbox}
        \end{column}
    \end{columns}

    \vspace{0.24em}

    \begin{itemize}\scriptsize
        \setlength\itemsep{0.1em}
        \item Kumar et al.\ (2020) argue that the mathematical appeal of Shapley values does not
              transfer to the human goals people bring to explanations, and that closing the gap
              requires causal assumptions the method does not make.
        \item If the user needs \textbf{recourse} (``what do I change?''), the right tool is a
              \textbf{counterfactual explanation} --- the nearest feasible input that flips the
              decision --- not a feature attribution.
        \item \textbf{Not unique, and not a defence.} Different methods, backgrounds and seeds give
              different attributions for the same prediction; a reviewer finding an explanation
              satisfying is not evidence the model is correct.
        \item \textbf{Rule of thumb:} use explanations to \textbf{generate hypotheses}, then test
              each one by intervention --- an ablation, a re-fit without the feature, a fresh slice.
    \end{itemize}

    \vspace{0.09em}
    {\scriptsize I.~E.~Kumar, S.~Venkatasubramanian, C.~Scheidegger and S.~Friedler, ``Problems with
    Shapley-value-based explanations as feature importance measures,'' ICML 2020,
    \href{https://arxiv.org/abs/2002.11097v2}{arXiv:2002.11097v2}.}
\end{frame}
